%!TEX root = ./../../springer-mono-format/book.tex
%!TEX root = ./../../thesis-trier.tex

\IN{synchronizing!group}
\IN{section-regular partition}
\IN{trivial section-regular partitions}

Let $G \le_G \mathcal S_n$.
We call $G$ \emph{synchronizing}~\cite{ArnoldS06} 
if it synchronizes every non-permutation. 
A partition $P_1, \ldots, P_m$ of $[n]$ is called
\emph{section-regular}~\cite{Neumann09} 
for~$G$ if there exists a subset $S\subseteq [n]$
that is a section for each partition of the form
 $\applyfunc{P_1}{g}, \ldots, \applyfunc{P_m}{g}$
where $g \in G$. Equivalently
for each $g \in G$ the set $\applyfunc{S}{g}$ 
is a section of $P_1, \ldots, P_m$.

\begin{theorem}[{Neumann~\cite[Note 6.1]{Neumann09}}]
\label{sync_max:thm:sync_equals_no_sect_regular}
	A permutation group $G \le_G \mathcal S_n$ is synchronizing
	if and only if there exists no non-trivial section-regular
	partition for $G$. 
\end{theorem}

Next, we show that synchronizing groups lie between the $2$-homogeneous and primitive groups.

\begin{proposition}
\label{sync_max:prop:sync-between}
\begin{enumerate}
\item \label{sync_max:prop:sync-between:1} (Arnold and Steinberg~\cite{ArnoldS06}) A synchronizing group is primitive.

\item \label{sync_max:prop:sync-between:2} Every $2$-homogeneous group 
is synchronizing.

\end{enumerate}
\end{proposition}
\begin{proof}
\ref{sync_max:prop:sync-between:1}:	This follows since a system of blocks is a section-regular partition
(for every section).

\ref{sync_max:prop:sync-between:2}:  Let $t \colon [n] \to [n]$
 be a transformation such that $t(a) = t(b)$ for distinct $a,b \in [n]$.
 As $G$ is $2$-homogeneous, for distinct $c,d \in [n]$
 there exists $g \in G$
 such that $\applyfunc{\{c,d\}}{g} = \{a,b\}$.
 So,  $\concat{g}{t}$ maps $\{c,d\}$ to a singleton set.
 Thus, every pair of states can be mapped to a singleton set.
 By the transformation semigroup version
 of the pair criterion (\cref{prelim:thm:pair_sync_condition} and \cref{prelim:remarks_trans_semigrp_semiautomata})
 we find that every $2$-homogeneous group 
 is synchronizing.
\end{proof}

\begin{example}
	Let $n \ge 4$
	and  $G = \langle (0 ~ 1 ~ \ldots ~ n - 1) \rangle_G$.
    Then $G$ is synchronizing if and only if~$n$ is prime.
	For assume~$n$ is prime
	and suppose $P_1, \ldots, P_m$ is a section-regular partition
	witnessed by $S \subseteq [n]$.
	Let $i \in \{1,\ldots,m\}$.
	For $y \in P_i$ 
	set 
	\[
	 R_y = \{ g \in G \mid \exists s \in S \colon \applyfunc{s}{g} = y \}.
	\]
	Then $G = \bigcup_{y \in P_i} R_y$
	and $R_y \cap R_{y'} = \emptyset$
	for distinct $y, y' \in P_i$.
	As $G$ is a regular permutation group, for $s \in S$ and $y \in P_i$
	there exists precisely one $g \in G$
	with $g(s) = y$.
	So, $|R_y| = |S|$, which
	gives $n = |G| = |S||P_i|$.
	As~$n$ is prime, $|P_i| \in \{1,n\}$.
	Hence, the partition is trivial.
	By \cref{sync_max:thm:sync_equals_no_sect_regular}, $G$ is synchronizing.
	Conversely, assume $n = rm$ for $0 < r,m < n$.
	Then $\Delta = \{ a \in [n] \mid a \equiv 0 \pmod{r} \}$
	is a non-trivial block.
	So, by \cref{sync_max:prop:sync-between:1} of \cref{sync_max:prop:sync-between}, $G$ 
	is not synchronizing.
\end{example}

Ara\'{u}jo asked whether the synchronizing groups are properly contained
inside the primitive groups (see the introduction of~\cite{Neumann09}).
This was answered affirmatively
by Neumann~\cite{Neumann09}.
The next example giving a non-synchronizing primitive group is due to Neumann~\cite[Ex.~3.3]{Neumann09}.

\begin{example}[{Neumann~\cite[Ex.~3.3]{Neumann09}}]
\label{sync_max:ex:prim_not_sync}
 Set $G = \mathcal A_5 \times \mathcal A_5$.
 It is well-known~\cite{cameron_1999,dixon1996permutation}
 that~$\mathcal A_5$ is simple and $Z(\mathcal A_5) = \{ \operatorname{id}_{[5]} \}$
 (recall that 
 $Z(G) = \{ g \in G \mid \forall h \in G \colon hg = gh \}$
 for a group~$G$, see page~\pageref{def:center}).
 By \cref{prelim:prop:diagonal_action},
 the representation of~$G$ by the diagonal action $G^D$
 as a permutation group on~$\mathcal A_5$ is primitive.
 Let~ $U$ be the point stabilizer subgroup of~$5$ in~$\mathcal A_5$, which
 is isomorphic to $\mathcal A_4$.
 Set $g = (0 ~ 1 ~ 2 ~ 3 ~ 4)$.
 
 Let $y \in G$. By \cref{prelim:eqn:stab}, $U^y$ is the point stabilizer of $y(5)$.
 By \cref{prelim:coset-point-stab}, 
 the sets $g^i U^y$ for $i \in \{0,1,2,3,4\}$ are all the left cosets of $U^y$.
 Thus, the set $S = \{\operatorname{id}_{[5]}, g, g^2, g^3, g^4 \}$
 is a section for the partition induced by all these cosets.
 Since for $i \in \{0,1,2,3,4\}$ and $(x,y) \in G$ we have
 $xg^iUy^{-1} = xy^{-1} (g^y)^i U^y$, the sets
 \[
 x U y^{-1}, x gU y^{-1}, x g^2 U y^{-1}, x g^3 U y^{-1}, x g^4 U y^{-1}
 \]
 are all the left cosets of $U^y$.
 So, the partition into left cosets of~$U$
 is section-regular for $G^D$. Thus, $G^D$ is not synchronizing.
\end{example}

McAlister~\cite{mcalister1998semigroups} (see also \cite[Thm. 8.2]{ArCamSt2017})
proved that a permutation group $G$
is $2$-homogenous if and only 
if $\langle G \cup \{t\}\rangle_S$
contains every non-permutation
for any $t \colon [n]\to [n]$ of defect one.
Next, we give a different characterization for $k$-homogeneity in \cref{lem:k_hom_char}.
This characterization was given in~\cite[Lem.~4.1]{HoffmannLATAextended2}.
Then, in \cref{thm:all-n-minus-k-sets-reachable}, 
we characterize those groups on domains of size~$n$
in which all subsets of size $n - i\cdot k$ are reachable
when adjoining a transformation of defect~$k$.
As a sufficient condition for subset reachability in the special case $k = 2$
this was proved in~\cite[Prop.~4.2]{HoffmannLATAextended2}.

\begin{proposition} \label{lem:k_hom_char}
	Let $G \le_G \mathcal S_n$ and $1 \le k \le n - 1$. Then the following are equivalent.
	\begin{enumerate}
		\item \label{cond:hom} $G$ is $k$-homogeneous.
		
		\item \label{cond:any_f} Every $(n-k)$-set is reachable in $\langle G \cup \{t\} \rangle_S$ for every $t \colon [n] \to [n]$ of defect $k$.
		
		\item \label{cond:some_f} Every $(n-k)$-set	is reachable in $\langle G \cup \{t\} \rangle_S$ for some $t \colon [n] \to [n]$ of defect $k$.
		
		\item \label{cond:any_f_k} Every $k$-set is reachable in $\langle G \cup \{t\} \rangle_S$ for every $t \colon [n] \to [n]$ of rank $k$.
		
		\item \label{cond:some_f_k} Every $k$-set is reachable in $\langle G \cup \{t\} \rangle_S$ for some $t \colon [n] \to [n]$ of rank $k$.
	\end{enumerate}		
\end{proposition}
\begin{proof}	
	We first show that \cref{cond:hom,cond:any_f,cond:some_f} are equivalent.
	Since $k$-homogeneity is equivalent to $(n-k)$-homogeneity, and $t$ has defect~$n-k$
	if and only if it has rank~$k$, it then follows that
	\cref{cond:hom,cond:any_f_k,cond:some_f_k} are also equivalent.
	
	$\ref{cond:hom} \Rightarrow \ref{cond:any_f}$:
	Let $A \subseteq [n]$ be of size $n-k$.
	Since $G$ is $k$-homogeneous, there exists $g \in G$ such that
	$
	[n] \setminus A = g\bigl([n] \setminus t([n])\bigr).
	$
	Consequently,
	$
	A = g(t([n])).
	$
	
	$\ref{cond:any_f} \Rightarrow \ref{cond:some_f}$:
	This is clear.
	
	$\ref{cond:some_f} \Rightarrow \ref{cond:hom}$:
	Let $A \subseteq [n]$ be of size $n-k$.
	By assumption, and since $|A| < n$, there exist
	$h \in \langle G \cup \{t\} \rangle_S$ and $g \in G$ such that
	\[
	A = g(t(h([n]))).
	\]
	Hence,
	$
	|t(h([n]))| = n-k.
	$
	Since $t(h([n])) \subseteq t([n])$ and $|t([n])| = n-k$, we must have
	$
	t(h([n])) = t([n]).
	$
	Therefore, every $(n-k)$-subset of $[n]$ is the image of $t([n])$ under an element
	of~$G$. Thus $G$ is $(n-k)$-homogeneous, and hence $k$-homogeneous.
\end{proof}

\begin{corollary} \label{cor:k_hom_char}
	Let $G \le_G \mathcal S_n$ be a permutation group and $2 \le k < n$.	
	Suppose for each  $t \colon [n] \to [n]$ of defect $k$,
	all subsets of size at most $n - k$ are reachable
	in $\langle G \cup \{t\} \rangle_S$.
	Then $G$ is $k$-homogeneous.
\end{corollary}

\begin{theorem}\label{thm:all-n-minus-k-sets-reachable}
	Let $G \le_G \mathcal S_n$ be a permutation group and let $1 \le k < n$.
	Then the following are equivalent:
	\begin{enumerate}
		\item
		For every [idempotent] one-point singular
		$t \colon [n] \to [n]$ of defect~$k$,
		all subsets of size $n-ik$ for
		$1 \le i < \frac{n}{k}$ are reachable in
		$\langle G \cup \{t\} \rangle_S$.
		
		\item
		For distinct
		$a_1,\ldots,a_{k+1}\in[n]$
		and every subset $A\subseteq[n]$ of size
		$n-ik$ for $1 \le i < \frac{n}{k}$,
		there exists $g\in G$ such that
		$a_1\in g(A)$ and
		$\{a_2,\ldots,a_{k+1}\}\cap g(A)=\emptyset$.
	\end{enumerate}
\end{theorem}
\begin{proof}	
	Let $t \colon [n] \to [n]$ be a one-point singular of defect $k$.
	Consider $A \subseteq [n]$ with $|A| = n - ik$ for some $1 \le i < \frac{n}{k}$.
	We proceed by induction on~$i$.
		
	Suppose first that $i = 1$.
	Since $\excl{t}$ is a $k$-set,
	by assumption, we have $\excl{t} \cap g(A) = \emptyset$ for some $g \in G$.
	As $|A| = n - k$, this gives $g(A) = t([n])$.
	Thus $A = g^{-1}(t([n]))$.
	
	Next, suppose $i > 1$. Let $\dupl{t} = \{c\}$.
	By assumption, there exists $g \in G$ such that
	\[
	c \in g(A)
	\quad\text{and}\quad
	\excl{t} \cap g(A) = \emptyset.
	\]
	Observe the following.
    
    \begin{enumerate}
    \item By \cref{prop:char-one-attractor}, since~$t$ is a one-point singular of defect~$k$,
    we have $c = t(a_1) = \cdots = t(a_{k})$ for $k$ distinct points $a_1, \ldots, a_k \in [n]$.
      
    \item 
    Since $g(A)$ contains no point of $\excl{t}$,
    every point of $g(A)$ lies in $t([n])$.
    Hence, for every $a \in A$ with $g(a) \neq c$,
    there exists $b \in [n]$ such that
    $g(a) = t(b)$.
    \end{enumerate}
    Both conditions yield that $B = t^{-1}(g(A))$ has size $|A| + k$
	and satisfies $t(B) = g(A)$.
	By the induction hypothesis, we have $B = h([n])$
	for some $h \in \langle G \cup \{t\} \rangle_S$.
	Hence, we find $A = g^{-1}(t(h([n])))$.		
	
	Conversely, let $a_1, \ldots, a_{k+1} \in [n]$ be distinct and 
	$A \subseteq [n]$ with $|A| = n - ik$ for some $1 \le i < \frac{n}{k}$.
	Consider the idempotent one-point singular of defect~$k$
	such that
	\[
	t(a_1)=t(a_2)=\cdots=t(a_{k+1})=a_1.
	\]
	Since~$A$ is reachable in $\langle G \cup \{t\} \rangle_S$
	and $A \ne [n]$,
	we have $A = g(t(h([n])))$
	for some $g \in G$ and $h \in \langle G \cup \{t\} \rangle_S$.
 	Since a permutation maps $[n]$ to itself, 	
 	we may assume $h = g_1 t \cdots t g_m t$
	for some $g_1, \ldots, g_m \in G$, where $m \ge 0$ is minimal.	
	We have $a_1 \in t(h([n]))$, since otherwise, none
	of the points $a_1, \ldots, a_{k+1}$ is in $h([n])$
	and so $t$ acts as the identity on the remaining points.
	In this case, $A = g(h([n]))$, contradicting the minimal
	choice of~$m$.
	Since $\excl{t} = \{a_2, \ldots, a_{k+1}\}$,
	none of these points belongs to $t(h([n]))$.
	Therefore,
	\[
	a_1\in g^{-1}(A)
	\quad\text{and}\quad
	\{a_2,\ldots,a_{k+1}\}\cap g^{-1}(A)=\emptyset,
	\]
	as required.
\end{proof}

Note that for $k=1$, \cref{thm:all-n-minus-k-sets-reachable}
reduces to \cref{sync_max:thm:prim_char:partition}
of \cref{sync_max:thm:prim_orbitals}.